% !TeX root = ../example.tex
\section{媒体与收尾}
\sectioncoverpage[assets/sigs_night]{本节演示视频链接、参考文献与致谢。}
\note{五个小节用于验证奇数项双栏布局。最后一项位于左列，随后演示媒体、文献与收尾。}
\subsection{视频链接}

\begin{frame}{视频链接示例}
      \centering
      % Match the demo video's 16:9 ratio so its first frame fills the area.
      \href{run:../assets/presentation_demo.mp4}{% 原创 CC0 动画；路径相对于 build/ 中的 PDF
            \XeTeXLinkBox{\parbox[c][.60\textheight][c]{\dimexpr .60\textheight*16/9\relax}{\mbox{}}}%
      }
      \campusvideocaption{原创动画：大纲 \(\rightarrow\) 幻灯片 \(\rightarrow\) 备注（6 秒，无音轨，CC0）。}
\end{frame}
\note{本页的原创几何动画展示大纲如何对应幻灯片及逐页备注，采用 CC0-1.0，无外部画面或音轨。PDF 中点击会打开本地视频；PPTX 内嵌视频和首帧预览，点击视频播放。}

\subsection{参考文献与致谢}

\begin{frame}[fragile]{参考文献与致谢用法}
      \begin{description}[汇总单页]
            \item[汇总单页] \texttt{\string\referencespage[shrink][\string\small]} 可指定字号并缩放到一页。
            \item[汇总续页] \texttt{\string\referencespage[break][\string\footnotesize]} 内容过多时自动续页。
            \item[双栏汇总] 命令末尾加 \texttt{[twocolumn]}；省略时默认为单栏。
            \item[致谢页面] \texttt{\string\backmatter} 会同步显示封面的标题、副标题、课程/课题组、作者、指导教师和日期；加入 \texttt{[notitle]} 可隐藏主标题。
      \end{description}
\end{frame}
\note{介绍参考文献的单页、续页与双栏选项。致谢页同步封面元信息，可选择隐藏主标题。}

\subsection{收尾说明}
\begin{frame}{收尾页选项}
  \begin{itemize}
    \item 参考文献页可以选择字号、续页和栏数。
    \item 致谢页可以保留或隐藏主标题。
  \end{itemize}
\end{frame}
\note{根据文献数量选择可读的参考文献布局，随后进入文献汇总与致谢。}

\subsection{参考文献汇总}
\referencespage[break][\small][onecolumn] % 第三个参数改为 [twocolumn] 可使用双栏
\note{这里汇总本演示引用的条目，按首次引用顺序编号。示例数据库的元数据仍需在正式使用前核实补全。}

\subsection{致谢}
\backmatter     % 不显示主标题时改为 \backmatter[notitle]
\note{结束演示，邀请反馈代码窗口、页面可读性和模板使用方式。}
